% ============================================================
% DOCUMENTCLASS
% ============================================================
\documentclass[12pt,oneside]{book}

% ============================================================
% METADATOS
% ============================================================
\newcommand{\universidad}{Universidad de Buenos Aires (UBA)}
\newcommand{\materia}{Derecho Procesal Civil}
\newcommand{\titulo}{Unidad 1 --- Introducción al Proceso Judicial}
\newcommand{\autor}{Isaac Edgar Camacho Ocampo}
\newcommand{\anio}{2025}

% ============================================================
% PAQUETES
% ============================================================
\usepackage[utf8]{inputenc}
\usepackage[T1]{fontenc}
\usepackage[spanish,es-nodecimaldot]{babel}

\usepackage[
    top=2.5cm,
    bottom=2.5cm,
    left=3cm,
    right=2.5cm,
    headheight=15pt
]{geometry}

\usepackage{amsmath}
\usepackage{amssymb}
\usepackage{mathtools}

\usepackage{graphicx}
\usepackage{float}
\usepackage{caption}
\usepackage{subcaption}

\usepackage{booktabs}
\usepackage{array}
\usepackage{longtable}
\usepackage{tabularx}

\usepackage{enumitem}

\usepackage{xcolor}
\usepackage[most]{tcolorbox}

\usepackage{fancyhdr}
\usepackage{ragged2e}

\usepackage[
    colorlinks=true,
    linkcolor=blue!50!black,
    citecolor=green!40!black,
    urlcolor=blue!60!black,
    pdftitle={\titulo},
    pdfauthor={\autor},
    pdfsubject={Apuntes universitarios},
    bookmarks=true
]{hyperref}

% ============================================================
% CONFIGURACIÓN
% ============================================================
\graphicspath{
    {./img/}
    {./graficos/}
    {./figuras/}
}

\setcounter{tocdepth}{2}

\setlength{\parindent}{0pt}
\setlength{\parskip}{0.5em}

\setlist[itemize]{
    topsep=0.3em,
    itemsep=0.2em
}

\setlist[enumerate]{
    topsep=0.3em,
    itemsep=0.2em
}

\clubpenalty=10000
\widowpenalty=10000

% ============================================================
% COMANDOS
% ============================================================
\newcommand{\abs}[1]{\left|#1\right|}
\newcommand{\norm}[1]{\left\lVert#1\right\rVert}
\newcommand{\vect}[1]{\mathbf{#1}}
\newcommand{\deriv}[2]{\frac{d#1}{d#2}}
\newcommand{\pderiv}[2]{\frac{\partial#1}{\partial#2}}

% ============================================================
% ENTORNOS / CAJAS
% ============================================================
\newtcolorbox{definition}[1]{
    enhanced,
    breakable,
    title={Definición: #1},
    fonttitle=\bfseries,
    colback=blue!3,
    colframe=blue!50!black,
    boxrule=0.8pt,
    arc=2mm
}

\newtcolorbox{important}{
    enhanced,
    breakable,
    title={Importante},
    fonttitle=\bfseries,
    colback=yellow!8,
    colframe=orange!70!black,
    boxrule=0.8pt,
    arc=2mm
}

\newtcolorbox{example}[1]{
    enhanced,
    breakable,
    title={Ejemplo: #1},
    fonttitle=\bfseries,
    colback=green!4,
    colframe=green!40!black,
    boxrule=0.8pt,
    arc=2mm
}

\newtcolorbox{formula}[1]{
    enhanced,
    breakable,
    title={Fórmula / Regla: #1},
    fonttitle=\bfseries,
    colback=purple!4,
    colframe=purple!50!black,
    boxrule=0.8pt,
    arc=2mm
}

\newtcolorbox{exercise}[1]{
    enhanced,
    breakable,
    title={Ejercicio: #1},
    fonttitle=\bfseries,
    colback=gray!5,
    colframe=gray!60!black,
    boxrule=0.8pt,
    arc=2mm
}

\newtcolorbox{note}{
    enhanced,
    breakable,
    title={Nota},
    fonttitle=\bfseries,
    colback=white,
    colframe=gray!50,
    boxrule=0.6pt,
    arc=2mm
}

% ============================================================
% CONFIGURACIÓN FINAL (encabezados y pies)
% ============================================================
\pagestyle{fancy}
\fancyhf{}

\fancyhead[L]{\nouppercase{\leftmark}}
\fancyhead[R]{\materia}

\fancyfoot[C]{\thepage}

\renewcommand{\headrulewidth}{0.4pt}
\renewcommand{\footrulewidth}{0pt}

\fancypagestyle{plain}{
    \fancyhf{}
    \fancyfoot[C]{\thepage}
    \renewcommand{\headrulewidth}{0pt}
}

% ============================================================
% DOCUMENT
% ============================================================
\begin{document}

% ---------------- PORTADA ----------------
\begin{titlepage}

\thispagestyle{empty}

\begin{center}

\vspace*{1cm}

{\large \textsc{\universidad}}

\vspace{3cm}

{\Huge \textbf{\materia}}

\vspace{1cm}

{\LARGE \titulo}

\vspace{0.8cm}

\textit{Comprender el proceso es el primer paso para transformar el estudio en criterio profesional.}

\vspace{1.2cm}

\rule{0.70\textwidth}{0.5pt}

\vspace{0.5cm}

{\large Apuntes de estudio}

\vspace{0.5cm}

\rule{0.70\textwidth}{0.5pt}

\vfill

{\large \autor}

\vspace{0.3cm}

{\large \anio}

\end{center}

\vfill

\begin{flushleft}
\textit{Este es un modesto aporte a los estudiantes de ``Derecho Procesal Civil'' no pretende sustituir el material oficial de la materia.}
\end{flushleft}

\end{titlepage}

% ---------------- ÍNDICE ----------------
\tableofcontents

% ============================================================
% CONTENIDO
% ============================================================

\chapter{Introducción al Proceso Judicial}

Esta unidad introductoria aborda los fundamentos del Derecho Procesal Civil, desde la función del abogado en la sociedad hasta los elementos técnicos del proceso judicial. Se desarrollan los conceptos de conflicto humano, proceso, sentencia, jurisdicción, competencia, cargas procesales, resoluciones judiciales y nulidad procesal, y se analiza la estructura general del Código Procesal Civil y Comercial de la Nación.

El conocimiento de las reglas y de la lógica del proceso judicial resulta relevante para el ejercicio profesional con independencia de la especialidad que se adopte, en la medida en que permite identificar al juez competente antes de presentar una demanda, distinguir entre cargas y obligaciones procesales, redactar escritos adecuados, detectar resoluciones nulas y sustentar su impugnación, y argumentar con solidez ante el tribunal.

\section{El abogado y el conflicto humano}

\subsection{La función del abogado en la sociedad}

La actividad profesional del abogado suele describirse, en un primer acercamiento, en relación con la defensa de los intereses de su cliente. Sin embargo, la función del abogado no se agota en la defensa de intereses particulares: comprende también el resguardo de las \textbf{garantías constitucionales} dentro del \textbf{debido proceso}. Estas garantías tienen como finalidad última asegurar la \textbf{igualdad} entre las partes frente a la resolución de los conflictos.

\subsection{El conflicto humano como materia de trabajo}

El \textbf{conflicto humano} constituye la materia prima sobre la que trabaja el abogado. Los conflictos son inherentes a la naturaleza humana, en tanto las personas poseen necesidades, emociones e intereses distintos que no siempre resultan conciliables por sí mismas, lo cual explica la necesidad de la intervención de un tercero imparcial para su resolución.

Quienes acuden a un estudio jurídico lo hacen, por regla general, motivados por una situación problemática (un despido laboral, el cierre de una empresa, un conflicto vinculado con el régimen de contacto con los hijos, entre otros supuestos), lo cual vincula la práctica profesional de manera constante con el conflicto. Este puede ser de naturaleza económica, familiar o personal, y detrás de cada expediente existe una persona que atraviesa una situación real, dimensión que no siempre es advertida por quienes intervienen en el proceso desde una posición de distancia burocrática.

\begin{example}{Efecto de la demora judicial}
Se plantea el caso de una persona que, tras el fallecimiento accidental de su cónyuge, debe atravesar un proceso judicial de aproximadamente diez años para obtener una indemnización, aun cuando esta resultara cuantiosa (del orden de cincuenta millones de pesos). El ejemplo ilustra que el prolongado estado de incertidumbre procesal genera un perjuicio autónomo, independiente del resultado económico final: la resolución del conflicto --- aun cuando no sea la más justa posible --- permite a la persona afectada continuar con su vida, mientras que la incertidumbre prolongada impide todo avance.
\end{example}

\subsection{La ética profesional del abogado}

La formación profesional del abogado debe orientarse a representar adecuadamente los intereses de su asistido sin recurrir a la fabricación o manipulación de pruebas ni a la comisión de actos ilegales. La responsabilidad de actuar conforme a la ética profesional recae exclusivamente en cada profesional, dado que ninguna circunstancia externa obliga a apartarse de la conducta correcta.

\section{El proceso judicial: concepto y estructura}

\subsection{Concepto de proceso judicial}

El término ``proceso'' remite a una secuencia de etapas que se inicia en un punto determinado y culmina en un resultado final. El \textbf{proceso judicial} constituye una especie particular de proceso, caracterizado por un alto grado de formalidad: se encuentra regulado por normas que establecen plazos y requisitos específicos, cuyo incumplimiento determina la ineficacia del acto procesal correspondiente.

\subsection{El fin del proceso: la sentencia y el ideal de justicia}

El fin del proceso judicial es la \textbf{sentencia}. Idealmente, dicha sentencia debería ser justa; sin embargo, la justicia constituye un ideal que no siempre resulta plenamente alcanzable, en la medida en que lo que resulta justo para una parte puede no serlo para la otra. La justicia opera, en tal sentido, como un horizonte hacia el cual el proceso judicial tiende, sin que ello implique la posibilidad de alcanzarla de manera exacta y absoluta en todos los casos.

\subsection{El debido proceso y la igualdad de condiciones}

Si la justicia perfecta no puede garantizarse de manera absoluta, el ordenamiento procesal sí puede garantizar que las reglas del proceso sean las mismas para todas las partes. El juez posee tanto la facultad como la \textbf{obligación} de resolver el conflicto sometido a su conocimiento, y el conjunto de normas procesales debe asegurar que ambas partes se encuentren en igualdad de condiciones frente a él.

\begin{definition}{Debido proceso}
Conjunto de normas procesales que garantizan que las dos partes de un conflicto judicial se encuentren en igualdad de condiciones frente al juez. Constituye la herramienta mediante la cual el ordenamiento jurídico procura aproximarse al ideal de justicia.
\end{definition}

\section{Estructura del Código Procesal Civil y Comercial de la Nación}

\subsection{Parte general y parte especial}

El Código Procesal Civil y Comercial de la Nación se organiza en una \textbf{parte general} y una \textbf{parte especial}. La parte general resulta aplicable a la generalidad del proceso y se encuentra dividida en distintos libros. Su desarrollo comienza con la regulación del órgano judicial, en particular con la figura del juez.

\subsection{El juez: jurisdicción e imperio}

El juez es definido como un tercero imparcial que adopta una decisión sobre la base de las posiciones sostenidas por cada una de las partes. A diferencia de otros funcionarios públicos, el juez posee \textbf{jurisdicción}.

\begin{definition}{Jurisdicción}
Facultad de decir el derecho (\emph{jus dicere}) y de ejecutarlo mediante el \textbf{imperio}, esto es, mediante la posibilidad de emplear la fuerza pública para hacer cumplir lo resuelto. Es una facultad exclusiva del juez: ningún otro funcionario público reúne ambas potestades.
\end{definition}

\begin{example}{Ejecución forzada de una sentencia civil}
Ante el incumplimiento de una obligación de pago --- por ejemplo, una deuda dineraria entre particulares ---, la parte acreedora puede iniciar un juicio que concluya en una sentencia condenatoria, notificada mediante cédula, que fije un plazo (por ejemplo, diez días) para el pago de la suma adeudada con más sus intereses. Vencido dicho plazo sin que se verifique el pago, el juez se encuentra facultado para disponer la ejecución forzada de la sentencia, incluyendo la venta de bienes del deudor a fin de satisfacer el crédito. Esta capacidad de emplear la fuerza para hacer cumplir lo resuelto es lo que distingue al juez de los demás funcionarios públicos.
\end{example}

En el ámbito del derecho administrativo, toda decisión de la administración pública admite revisión judicial; de no ser así, se estaría reconociendo jurisdicción a los funcionarios administrativos. Por ello, la garantía consiste en que la ejecución final de cualquier decisión requiera la intervención de un juez.

\begin{example}{Juicio de apremio}
Cuando la Administración Federal de Ingresos Públicos (AFIP) pretende cobrar impuestos o aplicar multas, no puede proceder a su ejecución directa, sino que debe recurrir ante un juez, único habilitado para disponer la ejecución. Dicho proceso se denomina \textbf{juicio de apremio}.
\end{example}

\subsection{La competencia: concepto, improrrogabilidad y prórroga}

La \textbf{competencia} puede definirse como un límite a la jurisdicción: el conjunto de asuntos y los límites dentro de los cuales un juez determinado se encuentra habilitado para resolver un conflicto concreto.

\begin{table}[H]
\centering
\caption{Criterios de competencia}
\begin{tabularx}{\textwidth}{>{\bfseries\raggedright\arraybackslash}p{0.28\textwidth} X}
\toprule
Criterio & Descripción \\
\midrule
Territorio & El lugar donde ocurrió el hecho o donde tiene domicilio alguna de las partes determina qué juez territorial es competente. \\
Fuero (materia) & La naturaleza del conflicto (civil, penal, comercial, laboral, etc.) determina qué fuero interviene. Un homicidio corresponde a un juez penal; una deuda, a un juez civil. \\
Improrrogabilidad (regla general) & La competencia es improrrogable: cada asunto tiene predeterminado un juez, y no se puede elegir libremente a cuál acudir. \\
Prórroga (excepción) & Solo cuando el conflicto es estrictamente patrimonial, sin cuestiones de orden público, de menores ni de familia, las partes pueden prorrogar la competencia a un tribunal distinto del que correspondería originalmente. \\
\bottomrule
\end{tabularx}
\end{table}

\begin{note}
Es habitual que, en contratos celebrados entre empresas, se incluya una cláusula que determine que los tribunales competentes en caso de conflicto serán los tribunales nacionales de la Capital Federal. Ello resulta admisible en la medida en que se trate de cuestiones estrictamente patrimoniales, respecto de las cuales la prórroga de competencia se encuentra permitida.
\end{note}

\subsection{Recusaciones, excusaciones, deberes y facultades del juez}

Con posterioridad a la regulación de la competencia, el Código regula las \textbf{recusaciones} y las \textbf{excusaciones}, institutos cuya finalidad es evitar que el juez actúe con parcialidad en la resolución del conflicto.

A continuación, el Código regula los \textbf{deberes y facultades del juez}, entre los que se incluyen el dictado de sentencias y de resoluciones intermedias. Por sobre estas funciones, el juez cumple un rol de director del proceso: debe velar por el cumplimiento de las normas procesales, las cuales tienen como finalidad constante asegurar la igualdad de condiciones entre las partes.

\section{Escritos, resoluciones judiciales y cargas procesales}

\subsection{Estructura del juzgado}

\begin{table}[H]
\centering
\caption{Jerarquía del juzgado}
\begin{tabularx}{\textwidth}{>{\bfseries\raggedright\arraybackslash}p{0.28\textwidth} X}
\toprule
Función & Descripción \\
\midrule
Mesa de entradas & Recibe y registra los escritos. Constituye la primera atención al público y el vínculo con los funcionarios del juzgado. \\
Jefe de despacho / Prosecretario & Puede dictar providencias simples, de mero trámite. No se encuentra habilitado para resolver controversias. \\
Secretario & Posee atribuciones para cuestiones de trámite de mayor complejidad, pero tampoco puede resolver controversias entre las partes. \\
Juez & Máxima autoridad del juzgado. Resuelve todas las controversias mediante resoluciones fundadas y dicta las sentencias definitivas e interlocutorias. \\
\bottomrule
\end{tabularx}
\end{table}

\subsection{El lenguaje de las partes y del juez}

Las partes se expresan en el proceso mediante \textbf{escritos}, cuyo contenido consiste en peticiones. El juez, por su parte, se expresa mediante \textbf{resoluciones}, empleando el modo imperativo (\emph{ordénase}, \emph{dígase}, \emph{hágase}), en tanto su función consiste en ordenar.

\subsection{Tipos de resoluciones judiciales (arts. 160 a 163 del Código Procesal)}

\begin{formula}{Clasificación de las resoluciones judiciales (arts. 160--163, Código Procesal Civil y Comercial de la Nación)}
Las resoluciones judiciales se clasifican en: (1) providencias simples; (2) sentencias interlocutorias; y (3) sentencias definitivas. Toda sentencia debe encontrarse razonablemente fundada.
\end{formula}

\begin{table}[H]
\centering
\caption{Tipos de resoluciones judiciales}
\begin{tabularx}{\textwidth}{>{\bfseries\raggedright\arraybackslash}p{0.22\textwidth} X}
\toprule
Tipo & Función / quién la dicta / exigencia de fundamentación \\
\midrule
Providencia simple & Mero trámite; impulsa el proceso sin resolver controversias (por ejemplo, librar un nuevo oficio). La dicta el jefe de despacho o el secretario. No requiere fundamentación extensa. \\
Sentencia interlocutoria & Resuelve una controversia surgida durante el proceso, por ejemplo, un incidente relativo a la competencia. Siempre la dicta el juez. Requiere sustanciación (traslado a la otra parte) y fundamentación. \\
Sentencia definitiva & Resuelve el fondo del asunto; es la decisión final del proceso. Siempre la dicta el juez. Requiere fundamentación completa y detallada. \\
\bottomrule
\end{tabularx}
\end{table}

Cuando una petición no genera controversia entre las partes --- por ejemplo, la solicitud de libramiento de un nuevo oficio ---, corresponde el dictado de una providencia simple. En cambio, cuando surge una controversia entre las partes --- por ejemplo, respecto de cuál es el juez competente para entender en la causa ---, dicha controversia debe resolverse mediante una sentencia interlocutoria.

\subsection{Cargas procesales vs. obligaciones}

\begin{important}
\textbf{Carga procesal:} conducta cuya omisión no genera sanción por parte del juez, pero coloca a quien no la realiza en una posición procesal desfavorable. Ejemplo: no ofrecer prueba testimonial no genera sanción, pero priva a la parte de la posibilidad de acreditar sus afirmaciones.

\textbf{Obligación civil:} conducta exigida cuyo incumplimiento genera una consecuencia jurídica directa, tal como la mora o los daños y perjuicios. Ejemplo: el incumplimiento en el pago del alquiler constituye en mora al inquilino y puede derivar en el desalojo.
\end{important}

Las partes no se encuentran obligadas a realizar ningún acto procesal en particular; sin embargo, el cumplimiento de las cargas procesales mejora su posición frente al juez. Por ejemplo, el ofrecimiento de prueba testimonial no es obligatorio, pero si dicha prueba resulta idónea para acreditar los hechos alegados, coloca a la parte que la ofrece en mejores condiciones para convencer al juez.

\subsection{La sustanciación de los incidentes}

La \textbf{sustanciación} consiste en el traslado de una petición formulada por una de las partes a la parte contraria, a fin de que esta última pueda expresar su posición antes de que el juez resuelva. De esta manera se garantiza que el juez cuente con ambas posiciones antes de adoptar una decisión, en resguardo del principio de contradicción y de la igualdad procesal.

\section{Caso práctico: nulidad procesal y fundamentación de las sentencias}

\subsection{Antecedentes del caso: la cámara Gessel}

\begin{example}{Nueva cámara Gessel solicitada faltando quince días para el juicio}
Se plantea un proceso penal en el cual el padre de un menor se encuentra acusado de abuso. Se habían realizado tres cámaras Gessel en el exterior, incorporadas al proceso mediante la vía diplomática, así como dos cámaras Gessel en el país que no resultaron útiles debido a que el menor prácticamente no había podido declarar. Estando ya fijada la fecha de juicio, la querella presentó certificados médicos y solicitó la realización de una nueva cámara Gessel, sosteniendo que el menor se encontraba en condiciones de declarar. Sin embargo, el período de ofrecimiento de prueba ya se encontraba cerrado al momento de dicha solicitud.
\end{example}

\begin{definition}{Preclusión}
Principio procesal según el cual, una vez vencida una etapa del proceso o el plazo previsto para la realización de un acto procesal, dicho acto ya no puede llevarse a cabo. Lo que ha precluido no puede reabrirse, salvo en los casos excepcionales expresamente previstos por la ley.
\end{definition}

\subsection{El trámite de sustanciación del incidente}

El pedido de la querella debía sustanciarse conforme al siguiente trámite:

\begin{enumerate}
    \item Vista al fiscal (parte acusadora pública).
    \item Vista al Ministerio Pupilar (por tratarse de un menor).
    \item Traslado a la defensa (representación del imputado).
\end{enumerate}

El fiscal se pronunció a favor de la realización de la nueva cámara Gessel, postura que fue acompañada por el Ministerio Pupilar. La defensa, por su parte, se opuso sobre la base de dos argumentos: en primer lugar, que la solicitud implicaba una dilación del proceso y que, habiendo operado la preclusión, no existía justificación para no haberla formulado con anterioridad, en particular al momento de incorporarse las traducciones de las cámaras realizadas en el exterior; en segundo lugar, que los certificados médicos acompañados por la querella no revestían el carácter de imparciales, por provenir de profesionales designados por la propia parte y no del cuerpo médico forense, que sería el organismo idóneo para efectuar dicha evaluación.

\subsection{La resolución del tribunal y la nulidad procesal}

El tribunal resolvió hacer lugar al pedido de la querella, remitiéndose a la conformidad prestada por la fiscalía y por el Ministerio Público, sin exponer fundamentos propios ni remitirse expresamente a los argumentos desarrollados por la fiscalía en su dictamen.

\begin{important}
La sola referencia a la conformidad de la querella y de la fiscalía no constituye, por sí misma, un fundamento jurídico válido. Para que el tribunal pudiera considerar propios los argumentos expuestos por la fiscalía, resultaba necesario que así lo manifestara expresamente, lo cual no ocurrió en el caso.
\end{important}

Frente a esta situación, se descarta la vía de la apelación, dado que las cuestiones relativas a la producción o denegación de prueba son, como regla, insusceptibles de dicho recurso; de admitirse la apelación de cada cuestión probatoria, la duración de los procesos se extendería de manera considerable. También se contempla la posibilidad de la revocatoria, calificada como un recurso de resultado incierto. Finalmente, se concluye que la vía adecuada es la \textbf{nulidad}, en tanto se encuentra comprometido el derecho de defensa.

\begin{definition}{Nulidad procesal}
Consecuencia que se produce cuando, en el curso del proceso, no se respeta una norma procesal o no se cumple con un requisito esencial. Se distingue de la situación en que el tribunal expone la totalidad de sus fundamentos: cuando, en cambio, la decisión se adopta de manera arbitraria, sin fundamento, o con un fundamento meramente aparente, el acto resulta pasible de ser tachado de nulo.
\end{definition}

\begin{formula}{Art. 163 CPCN --- fundamentación de las sentencias}
Toda sentencia debe encontrarse razonablemente fundada. Esta exigencia rige tanto para la sentencia definitiva como para las sentencias interlocutorias, con la única distinción de que las providencias simples, al no resolver controversias, no requieren una fundamentación extensa.
\end{formula}

Los fundamentos exigidos a una sentencia interlocutoria pueden ser de menor extensión que los requeridos para una sentencia definitiva, sin que ello implique una diferencia sustancial: en todos los casos debe existir un fundamento. La ausencia de fundamento, o la presencia de un fundamento meramente aparente, determina la nulidad de la resolución.

% ============================================================
% CORNELL
% ============================================================
\section*{Cuadro Cornell de repaso}
\addcontentsline{toc}{section}{Cuadro Cornell de repaso}

\begin{longtable}{
    >{\bfseries\raggedright\arraybackslash}p{0.30\textwidth}
    p{0.60\textwidth}
}
\caption{Cuadro Cornell --- Introducción al Proceso Judicial} \\
\toprule
\textbf{Preguntas / palabras clave} & \textbf{Notas / respuestas} \\
\midrule
\endfirsthead

\multicolumn{2}{c}{\textit{(continuación)}} \\
\toprule
\textbf{Preguntas / palabras clave} & \textbf{Notas / respuestas} \\
\midrule
\endhead

\bottomrule
\endfoot

\bottomrule
\endlastfoot

\textbf{¿Para qué sirve el abogado?} &
Su función no se limita a defender intereses de clientes. También comprende garantizar el respeto de las garantías constitucionales y el debido proceso, para que los conflictos se resuelvan en igualdad de condiciones. \\

\textbf{¿Qué es el conflicto humano?} &
Una pugna de intereses inherente a la naturaleza humana. Es inevitable porque las personas tienen necesidades, emociones e intereses distintos. Cuando no pueden resolverlo entre sí, requieren de un tercero (el juez). \\

\textbf{¿Qué es el proceso judicial?} &
Un conjunto estructurado de normas, dividido en etapas, que se inicia con la demanda y culmina con la sentencia. Es sumamente formal y puntilloso, a fin de garantizar la igualdad entre las partes. \\

\textbf{¿Se puede garantizar la justicia?} &
La justicia es un ideal. Como tal, constituye un camino, pero en sentido absoluto resulta inalcanzable. Sí se puede garantizar el debido proceso: que las reglas sean las mismas para todos. \\

\textbf{¿Qué es el debido proceso?} &
El conjunto de normas que garantiza que ambas partes en un litigio estén en igualdad de condiciones frente al juez. Es la herramienta que el ordenamiento jurídico ofrece para aproximarse al ideal de justicia. \\

\textbf{¿Qué es la jurisdicción?} &
La facultad que tiene exclusivamente el juez de decir el derecho (\emph{jus dicere}) y de ejecutarlo mediante el imperio (la fuerza). Ningún otro funcionario público puede reunir ambas potestades: solo el juez. \\

\textbf{¿Qué es la competencia?} &
El límite dentro del cual un juez puede ejercer su jurisdicción. Está determinada por el territorio y el fuero (materia). En principio es improrrogable; solo puede prorrogarse en cuestiones estrictamente patrimoniales. \\

\textbf{¿Qué son las resoluciones judiciales?} &
Las formas en que el juez provee los pedidos de las partes. Se dividen en providencias simples (trámite, sin controversia), sentencias interlocutorias (resuelven incidentes durante el proceso) y sentencias definitivas (resuelven el fondo). \\

\textbf{¿Qué es una carga procesal?} &
Una conducta procesal cuya realización coloca a la parte en mejor posición para lograr su objetivo (convencer al juez). Su omisión no genera sanción, pero sí una posición desventajosa. Se diferencia de la obligación, cuyo incumplimiento genera consecuencias jurídicas directas. \\

\textbf{¿Qué es la sustanciación?} &
El acto de trasladar una petición de una parte a la contraria para que esta pueda expresar su posición antes de que el juez resuelva. Garantiza el principio de contradicción y la igualdad procesal. \\

\textbf{¿Qué es la preclusión?} &
El principio por el cual, vencida una etapa o un plazo procesal, ya no puede realizarse el acto omitido. Garantiza el avance ordenado del proceso y evita que las partes reabran etapas ya superadas. \\

\textbf{¿Qué es la nulidad procesal?} &
La consecuencia que se produce cuando en el proceso se omite un requisito esencial establecido por la ley. Una resolución sin fundamentación puede ser tachada de nula por no cumplir con los recaudos del art. 163 CPCN. \\

\textbf{¿Por qué toda sentencia debe estar fundada?} &
Porque la fundamentación garantiza que la decisión no sea arbitraria, permite a las partes conocer las razones del fallo y ejercer su derecho de defensa. Sin fundamentación, la resolución puede impugnarse por nulidad. \\

\midrule
\multicolumn{2}{p{0.92\textwidth}}{
\textbf{Resumen:} Esta unidad introductoria de Derecho Procesal Civil establece las bases conceptuales de la materia. El abogado es, ante todo, un garante del debido proceso y de la igualdad de condiciones entre las partes, dado que la justicia perfecta es un ideal inalcanzable de manera absoluta. El proceso judicial es el mecanismo institucional que, a través de un conjunto riguroso de normas formales (el Código Procesal), intenta aproximarse a ese ideal. Dicho Código se organiza en una parte general --- que regula el órgano judicial, la jurisdicción, la competencia, las recusaciones, los deberes y facultades del juez, y la estructura del juzgado --- y una parte especial. El juez es el actor central del proceso: tiene jurisdicción (facultad de decir y aplicar el derecho) e imperio (facultad de ejecutar la decisión con la fuerza del Estado). La competencia limita esa jurisdicción y es improrrogable, salvo en cuestiones puramente patrimoniales. Los abogados se expresan mediante escritos (peticiones); el juez responde mediante resoluciones (providencias simples, sentencias interlocutorias o sentencias definitivas), siempre en tiempo imperativo, y toda resolución que resuelva una controversia debe estar razonablemente fundada, bajo pena de nulidad. Las partes actúan bajo cargas procesales, no obligaciones: omitir una carga no genera sanción, pero debilita la posición procesal. El caso práctico de la cámara Gessel ilustra estos conceptos en acción: preclusión, sustanciación, ausencia de fundamentación y nulidad procesal.
} \\

\end{longtable}

\end{document}
